\section{Shortcomings, Outstanding Items, and Lessons Learned}

\noindent Although our system is very robust to most scenarios there are several that it cannot handle. The following are the primary limitations of our system:

\begin{enumerate}
    \item The robot is not physically capable of separating $N\times M$ (with $N,M>2$) connections of blocks because it cannot grab any of the blocks to separate them.
    \item The robot may bump the structure in rare circumstances (but is able to recover from many of these accidents).
    \item The robot jerks suddenly if the motors briefly lose connection or the motor commands are sent slightly late.
    \item The robot cannot grasp some edge cases of misplaced (either by itself or human placed) blocks.
    \item The robot cannot "learn" from its memory/experiences of misgrasped and misplaced blocks.
\end{enumerate}

\noindent We have made several algorithm choices to alleviate the affect of these system limitations, but they may still arise. Firstly, the robot will ignore all blocks it has no way to grasp (e.g. blocks inside $N\times N$ connections or outside the workspace) and will only use blocks it can access. To minimize the chance of the robot bumping some part of the structure, we build the structure outwards and upwards as this allows to avoid placing blocks between or below others, which are the biggest causes of the robot damaging the structure. We have also considered (though didn't have time to implement) more significant fixes such as planning a trajectory using Probability Road Maps and A* or some other planning algorithms that would consider the structure as an obstacle. To soften the jerks, we implemented a "catching" behavior, which will recognize when there is a significant divergence in the robot's joint commands and true joint states and will "catch" the robot before it jerks by restarting the current trajectory with the initial position being the true robot joint states rather than forcing the robot to catch up to the trajectory by jerking towards joint states that it is supposed to be in. For misplaced blocks near/on the structure that the robot cannot grasp, it can detect these "unreachable" blocks and signal to the human that it needs help via a nod gesture. Finally, to alleviate the robot's inability to "learn" (or calibrate) kinematics+detection errors on the fly, we implemented an online calibration phase which produces an error map which we use in our ikin solves. 

\noindent Aside from the initial hardware constraint, the remaining limitations could theoretically be addressed by training and deploying a learned policy. We thoroughly considered this approach---specifically, using reinforcement learning in simulation with domain randomization to train a generalized grasping and placing policy. However, we ultimately concluded that the associated risks and costs (computational overhead, training feasibility, sim-to-real transfer challenges, and extended iteration times) outweighed the likely incremental benefits, especially considering we only had 7 weeks. Through rigorous system engineering, our current rules-based approach already provides a highly robust solution.
